\begin{manuallistedformatdiagram}{SAVE/RESTORE Base Header}{4}
\manualformatrowrange{header[63:0]}{63}{0}{%
\manualformatfield{reserved}{12}
\manualformatfield{STATUS}{16}
\manualformatfield{FLAGS}{16}
\manualformatfield{reserved}{11}
\manualformatfield{GSV}{5}
\manualformatfield{FMT}{4}
}
\manualformatrowrange{bitmap[63:0]}{63}{0}{%
\manualformatfield{state block bitmap}{64}
}
\end{manuallistedformatdiagram}
{\small\texttt{GSV} is the five-bit \texttt{GS\_VALID} bitmap; \texttt{FMT} selects the save-image format.\par}
\Needspace{2.65in}
\manualtablecaption{SAVE/RESTORE Fixed Base Block}
\begin{center}
\begin{tikzpicture}[x=1in,y=0.40in,every node/.style={font=\scriptsize}]
\def\manualsaveslot{0.63}
\node[anchor=west,font=\scriptsize\bfseries] at (0,0.34) {offset from the 4 KiB-aligned save-area base};
\node[anchor=east,text=ManualMuted] at ({8*\manualsaveslot},0.34) {offsets increase downward};

\node[anchor=east] at (-0.12,-0.40) {\texttt{0x000--0x00f}};
\path[manualLayoutHeader] (0,0) rectangle ({4*\manualsaveslot},-0.80);
\node[align=center] at ({2*\manualsaveslot},-0.40) {BASE HEADER\\\texttt{0x000} \textperiodcentered\ 8 bytes};
\path[manualLayoutHeader] ({4*\manualsaveslot},0) rectangle ({8*\manualsaveslot},-0.80);
\node[align=center] at ({6*\manualsaveslot},-0.40) {STATE BLOCK BITMAP\\\texttt{0x008} \textperiodcentered\ 8 bytes};

\node[anchor=east] at (-0.12,-1.20) {\texttt{0x010--0x04f}};
\foreach \reg in {0,...,7} {
  \path[manualLayoutSlot] ({\reg*\manualsaveslot},-0.80) rectangle ({(\reg+1)*\manualsaveslot},-1.60);
  \node at ({(\reg+0.5)*\manualsaveslot},-1.20) {\texttt{R\reg}};
}

\node[anchor=east] at (-0.12,-2.00) {\texttt{0x050--0x08f}};
\foreach \reg [evaluate=\reg as \reglabel using int(\reg+8)] in {0,...,7} {
  \path[manualLayoutSlot] ({\reg*\manualsaveslot},-1.60) rectangle ({(\reg+1)*\manualsaveslot},-2.40);
  \node at ({(\reg+0.5)*\manualsaveslot},-2.00) {\texttt{R\reglabel}};
}

\node[anchor=east] at (-0.12,-2.80) {\texttt{0x090--0x0bf}};
\foreach \seg in {0,...,4} {
  \path[manualLayoutOptional] ({\seg*\manualsaveslot},-2.40) rectangle ({(\seg+1)*\manualsaveslot},-3.20);
  \node at ({(\seg+0.5)*\manualsaveslot},-2.80) {\texttt{GS\seg}};
}
\path[manualLayoutReserved] ({5*\manualsaveslot},-2.40) rectangle ({6*\manualsaveslot},-3.20);
\node[align=center] at ({5.5*\manualsaveslot},-2.80) {reserved\\\texttt{0x0b8}};
\node[anchor=west,align=left,text=ManualMuted] at ({6.18*\manualsaveslot},-2.80) {GS images are optional\\and selected by \texttt{GS\_VALID}};

\node[anchor=east] at (-0.12,-3.72) {\texttt{0x0c0+}};
\path[manualLayoutExtension] (0,-3.20) rectangle ({8*\manualsaveslot},-4.24);
\node[align=center] at ({4*\manualsaveslot},-3.72) {EXTENSION STATE COMPONENTS\\CPUID-defined slots \textperiodcentered\ each component starts on a 64-byte boundary};

\node[anchor=west,text=ManualMuted] at (0,-4.55) {Each \texttt{Rn} and \texttt{GSn} cell is one 8-byte slot.};
\node[anchor=east,text=ManualMuted] at ({8*\manualsaveslot},-4.55) {\texttt{Rn}=\texttt{0x010}+8n; \texttt{GSn}=\texttt{0x090}+8n.};
\end{tikzpicture}
\end{center}
